现在可以陈述 Green 公式了.

\begin{Theorem}
\label{thm:elliptic-green-formula-adjoint-boundary-operators}
设 $L(\bm{x},D)$ 为 $\overline\Omega$ 上的 $2m$ 阶椭圆算子, $B_j(\bm{x},D)$, $j=1,\ldots,m$, 为一组正规边界算子. 设 $S_j(\bm{x},D)$, $j=1,\ldots,m$, 为一组边界算子, 将 $B_j$ 补全为 Dirichlet 系统. 则存在边界算子 $C_j(\bm{x},D)$, $T_j(\bm{x},D)$, $j=1,\ldots,m$, 具有以下性质:\begin{enumerate}\item $\operatorname{ord}C_j=2m-1-\operatorname{ord}S_j$, $\operatorname{ord}T_j=2m-1-\operatorname{ord}B_j$. ($\operatorname{ord}$ 表示算子的阶数.)\item $C_j$ 和 $T_j$ 构成 Dirichlet 系统.\item 对每个 $u,v\in H^{2m}(\Omega)$, 有
\begin{equation}\label{eq:elliptic-green-formula-adjoint-boundary-pairing}
\int_\Omega(Lu)v-u(L^*v)\,\mathrm{d}\bm{x}=\sum_{j=1}^m\int_{\partial\Omega}(S_ju)(C_jv)-(B_ju)(T_jv)\,\mathrm{d}S.
\end{equation}
这里 $L^*$ 为 $L$ 的形式伴随算子; 见 Definition~\ref{def:distribution-convolution-formal}.\item 若 $B_j$ 对 $L$ 满足补条件, 则 $C_j$ 对 $L^*$ 满足补条件.\end{enumerate}
\end{Theorem}

\begin{Proof}
分部积分得到如下形式的公式
\begin{equation}\label{eq:green-formula-raw-boundary-derivative-sum}
\int_\Omega(Lu)v-u(L^*v)\,\mathrm{d}\bm{x}=\sum_{\alpha,\beta}\int_{\partial\Omega}a_{\alpha\beta}(\bm{x})D^\alpha uD^\beta v\,\mathrm{d}S,
\end{equation}
其中对满足 $|\alpha|+|\beta|\leq2m-1$ 的 $\alpha$ 和 $\beta$ 求和. 接着在 $\partial\Omega$ 上分部积分, 将所有切向导数从 $u$ 转移到 $v$ 上, 使 $u$ 只剩纯法向导数 (这一步需要 $\partial\Omega$ 上的单位分解和局部坐标图). 由此得到如下形式的公式
\begin{equation}\label{eq:green-formula-pure-normal-boundary-reduction}
\int_\Omega(Lu)v-u(L^*v)\,\mathrm{d}\bm{x}=\sum_{j=0}^{2m-1}\int_{\partial\Omega}a_j(\bm{x})\frac{\partial^ju}{\partial n^j}E_j(\bm{x},D)v\,\mathrm{d}S,
\end{equation}
其中 $E_j$ 为 $2m-j-1$ 阶微分算子. 若 $L$ 椭圆 (甚至更一般地, 若 $\partial\Omega$ 非特征), 则 $E_j$ 含有与 $\dfrac{\partial^{2m-j-1}v}{\partial n^{2m-j-1}}$ 成比例且系数非零的项; 换言之, $E_j$ 构成 Dirichlet 系统. 接着利用 Lemma~\ref{lem:dirichlet-system-tangential-triangular-reduction} 得到
\begin{equation}\label{eq:green-formula-normal-derivative-dirichlet-system-expansion}
\frac{\partial^ju}{\partial n^j}=\sum_{k=1}^m\Psi_{jk}(\bm{x},D)B_k(\bm{x},D)u+\sum_{k=1}^m\Psi_{jk}^{\prime}(\bm{x},D)S_k(\bm{x},D)u.
\end{equation}
将其代入 \eqref{eq:green-formula-pure-normal-boundary-reduction}, 然后在 $\partial\Omega$ 上分部积分, 将切向微分算子 $\Psi_{jk}$ 和 $\Psi_{jk}^{\prime}$ 从 $u$ 转移到 $v$ 上. 这就得到 \eqref{eq:elliptic-green-formula-adjoint-boundary-pairing}.

为验证补条件, 设 $\Omega$ 为半空间 $\{x_n>0\}$, $L$ 具有常系数. 考虑形如 $\exp(i\bm{\xi}\cdot\bm{x})v(x_n)$ 的解, 其中 $\xi_n=0$, 并且当 $x_n\to\infty$ 时 $v(x_n)\to0$. 将 $\Omega$ 替换为 $\Sigma\times\mathbb{R}^+$ 后, Green 公式 \eqref{eq:elliptic-green-formula-adjoint-boundary-pairing} 仍成立, 其中 $\Sigma$ 是 $\mathbb{R}^{n-1}$ 中对应于一个周期的平行多面体. 此外, 此时 $L$ 成为常微分算子. 对这样的算子, Fredholm 二择一定理成立, 即初值问题
\begin{equation}\label{eq:green-complementing-homogeneous-normal-ode}
L(\bm{x},D)u=0,\ B_j(\bm{x},D)u=0\quad\text{for }x_n=0
\end{equation}
只有平凡解, 当且仅当问题
\begin{equation}\label{eq:green-complementing-inhomogeneous-normal-ode}
L(\bm{x},D)u=f,\ B_j(\bm{x},D)u=g_j\quad\text{for }x_n=0
\end{equation}
对所有 $f$ 和 $g_j$ 都可解. 由 Green 公式, 后一个条件意味着初值问题
\begin{equation}\label{eq:green-complementing-adjoint-homogeneous-normal-ode}
L^*(\bm{x},D)v=0,\ C_j(\bm{x},D)v=0\quad\text{for }x_n=0
\end{equation}
只有平凡解, 即 $C_j$ 对 $L^*$ 满足补条件. 注意, 若 $v$ 为 \eqref{eq:green-complementing-adjoint-homogeneous-normal-ode} 的非平凡解, 则 \eqref{eq:green-complementing-inhomogeneous-normal-ode} 中的 $f$ 和 $g_j$ 必须满足
\begin{equation}\label{eq:green-adjoint-mode-data-compatibility}
\int_\Omega fv\,\mathrm{d}\bm{x}=-\sum_{j=1}^m\int_{\partial\Omega}g_jT_jv\,\mathrm{d}\bm{x}.
\end{equation}
证明完成.
\end{Proof}
